\section{Relational probability: exact couplings with an optimality dual}
\label{sec:transport}

\subsection{Why independent execution is the wrong default}

Let two programs return fair bits. Their output distributions are equal, but
two \emph{independent} runs agree only with probability $1/2$. A proof of equality
in distribution must not require independent executions to return equal values.
Instead, it can exhibit a joint distribution with the required marginals that
puts both bits on the diagonal. Couplings are a standard foundation for such
relational proofs~\cite{coupling}. The new worker computes them exactly and
also explains when a perfect coupling is impossible.

\begin{definition}[Finite coupling and defect]
Let $\mu:X\to\Q_{\ge0}$ and $\nu:Y\to\Q_{\ge0}$ be probability distributions on
nonempty finite sets. A coupling is a matrix $\gamma:X\times Y\to\Q_{\ge0}$
with row sums $\mu$ and column sums $\nu$. For a relation $R\subseteq X\times Y$,
its defect is
\[
 \delta_R(\gamma)=\sum_{(x,y)\notin R}\gamma(x,y).
\]
A zero-defect coupling is supported by $R$.
\end{definition}

The problem supplied externally to the checker consists of the two complete
marginals and the relation. The certificate supplies the \emph{full} coupling,
not only its favorable entries. Otherwise a proposer could hide missing mass
and report a zero defect for two incompatible distributions.

\subsection{The Hall lower bound}

For $A\subseteq X$, put $R(A)=\{y:\exists x\in A,\ R(x,y)\}$ and define
\[
 h(A)=\mu(A)-\nu(R(A)).
\]
The empty subset gives $h(\varnothing)=0$.

\begin{lemma}[Weighted Hall obstruction]
\label{lem:hall}
For every coupling $\gamma$ and every $A\subseteq X$,
$\delta_R(\gamma)\ge h(A)$.
\end{lemma}
\begin{proof}
The mass in rows $A$ is $\mu(A)$. At most $\nu(R(A))$ of this mass can lie in
columns $R(A)$, by the column marginal equations and nonnegativity. Every entry
in $A\times(Y\setminus R(A))$ is outside $R$. Its total mass is therefore a
lower bound on the defect, proving the inequality.
\end{proof}

A certificate contains a coupling $\gamma$, a claimed defect $\delta$, and one
subset $A$. The checker verifies nonnegativity, both marginals, the actual sum
outside $R$, and
\begin{equation}
 \delta_R(\gamma)=\delta=h(A).
 \label{eq:transport-dual}
\end{equation}
No flow network is reconstructed by the checker.

\begin{theorem}[Optimality from a primal--dual equality]
\label{thm:transport}
If the transport certificate is accepted, $\gamma$ minimizes defect among all
couplings of $\mu$ and $\nu$, and $\delta$ is the exact minimum. In particular,
$\delta>0$ proves that no coupling supported by $R$ exists.
\end{theorem}
\begin{proof}
The supplied matrix attains $\delta$, while Lemma~\ref{lem:hall} gives the same
lower bound for every other coupling.
\end{proof}

This lower-bound evidence is stronger than ``the flow search failed to find a
perfect matching.'' It is also narrower than a claim about arbitrary source
programs: it proves an obstruction for these marginals and this relation. The
source bridge must establish that those are the ones in the original goal.

\subsection{An exact proposer}

Build a bipartite flow network with source $s$, sink $t$, an edge
$s\to x$ of capacity $\mu(x)$, an edge $y\to t$ of capacity $\nu(y)$, and
$x\to y$ of capacity 1 exactly when $R(x,y)$. A maximum flow has value $f$ and
produces a partial supported coupling $g$. The prototype uses Edmonds--Karp
augmenting paths with Python \code{Fraction}; there is no scaling to an
approximate integer grid.

Let residual row masses and column masses be
\[
 r_x=\mu(x)-\sum_y g(x,y),\qquad
 c_y=\nu(y)-\sum_x g(x,y),\qquad
 \delta=\sum_xr_x=\sum_yc_y=1-f.
\]
If $\delta=0$, take $\gamma=g$ and $A=\varnothing$. Otherwise complete the
partial coupling by
\begin{equation}
 \gamma(x,y)=g(x,y)+\frac{r_xc_y}{\delta}.
 \label{eq:residual-coupling}
\end{equation}
The added matrix has the missing marginals. Moreover, $r_xc_y>0$ implies
$\neg R(x,y)$: an admissible such pair would permit another positive
augmentation. Since $f<1$, its capacity-1 edge cannot already be saturated by
a total flow of size less than 1. Hence the completion adds exactly $\delta$
unsupported mass.

Let $A$ be the left vertices reachable from $s$ in the residual network. The
standard residual cut, written in terms of row and column masses, yields
$\mu(A)-\nu(R(A))=1-f$. The checker verifies this equality directly, so an
error in cut extraction is harmless to acceptance. This proves that every
finite rational transport input has a certificate of the specified form.

\begin{lstlisting}
proposeTransport(mu, nu, R):
  g, residual := exactMaxFlow(bipartiteNetwork(mu, nu, R))
  r := missing row masses of g
  c := missing column masses of g
  delta := sum(r)
  joint := g if delta = 0 else g + outer(r,c)/delta
  A := empty if delta = 0 else residualReachableLeftVertices
  return joint, delta, A

checkTransport(problem, certificate):
  check joint >= 0, rowSums(joint) = mu, columnSums(joint) = nu
  bad := sum of joint entries outside R
  check bad = claimedDelta = mu(A) - nu(R(A))
\end{lstlisting}

The dense checker takes $O(|X||Y|)$ rational arithmetic operations. The
Edmonds--Karp combinatorial bound is $O(VE^2)$ with
$V=|X|+|Y|+2$; the delivered dense residual representation is intended for
small instances, not high-throughput graph processing. Rational bit complexity
is additional. Production search may use a faster exact flow library or a
numerical proposal followed by exact repair, but acceptance remains
\eqref{eq:transport-dual}.

\subsection{Worked examples}

\paragraph{Monotone Bernoulli coupling.}
Let $X\sim\operatorname{Bernoulli}(1/3)$,
$Y\sim\operatorname{Bernoulli}(2/3)$, and let $R(x,y)$ mean $x\le y$. With rows
and columns ordered $0,1$, a perfect coupling is
\[
 \gamma=\begin{pmatrix}1/3&1/3\\0&1/3\end{pmatrix}.
\]
The independent product has positive mass at $(1,0)$, so it cannot prove the
same almost-sure relation. The distinction is the choice of a joint law, not
an improvement in numerical precision.

\paragraph{A proof that no perfect monotone coupling exists.}
Reverse the parameters. Rows have masses $(1/3,2/3)$, columns have masses
$(2/3,1/3)$, and the same order relation is requested. Taking $A=\{1\}$ gives
$R(A)=\{1\}$ and $h(A)=1/3$. The matrix
\[
 \gamma=\begin{pmatrix}1/3&0\\1/3&1/3\end{pmatrix}
\]
has defect exactly $1/3$, so its optimality is settled by three short sum
checks. A positive optimal defect is a successful answer, not a failed search.

\paragraph{Two coins versus one binomial sample.}
Let $X$ be the uniformly distributed pair $(a,b)\in\{0,1\}^2$, let
$Y\sim\operatorname{Binomial}(2,1/2)$, and require $a+b=Y$. Coupling each pair
with its sum gives defect zero. Under independent execution, agreement has
probability
\[
 (1/4)^2+(1/2)^2+(1/4)^2=3/8,
\]
so the defect is $5/8$. This is an explicit distributional-refinement target
that independent product execution cannot certify. A runnable example exports
both the witness and the independently computed defect.

\subsection{Composition with finite probabilistic programs}

For finite distributions, define
\[
 (\mu\mathbin{\gg=}K)(x')=\sum_x\mu(x)K(x,x').
\]
Suppose $\gamma$ couples $\mu$ and $\nu$, and for each pair $(x,y)$ of positive
joint mass, $\Gamma_{x,y}$ couples $K_x$ and $L_y$. Then
\begin{equation}
 \gamma'(x',y')=\sum_{x,y}\gamma(x,y)\Gamma_{x,y}(x',y')
 \label{eq:bind-coupling}
\end{equation}
is a coupling of the two bound distributions. The first marginal follows by
summing over $y'$, substituting the first marginal of each $\Gamma_{x,y}$,
and then substituting the first marginal of $\gamma$; the second is symmetric.
This finite-sum identity is the initial Lean composition theorem to prove.
It corresponds to the existing \code{PMF.bind} semantics, whose pointwise
formula and associativity law are present at the inspected Mathlib
pin~\cite{mathlib-pmf}.

For related state pairs, suppose each proposed transition coupling leaves the
relation with probability at most $\varepsilon\in[0,1]$. Starting in the
relation, the probability of leaving it within $k$ steps is at most
\[
 1-(1-\varepsilon)^k\le k\varepsilon.
\]
This follows by conditioning on the first $j$ steps all remaining related and
induction on $j$. Once outside the relation, any coupling of the marginals may
be used. The argument concerns a finite synchronous horizon; it does not claim
an asynchronous product-program theorem or a general infinite-path measure
construction. Those are valuable extensions, but a finite marginal checker
should not inherit their names without their proofs.

\subsection{A precise next extension: transport costs}

The defect problem is a special case of minimizing
$\sum_{x,y}c(x,y)\gamma(x,y)$, with cost $c=\one{\neg R}$. A general
rational-cost worker can retain the full coupling and replace the Hall subset
by potentials $u_x,v_y$ satisfying
\[
 u_x+v_y\le c(x,y),\qquad
 \sum_xu_x\mu(x)+\sum_yv_y\nu(y)
   =\sum_{x,y}c(x,y)\gamma(x,y).
\]
The inequality gives a lower bound after multiplication by $\gamma$ and
summation; equality proves optimality. This is a concrete certificate format
for min-cost transport, not an implemented feature in this archive.

With a distance cost, a family of witnesses satisfying
$\E[d(X',Y')]\le\kappa d(x,y)$ gives a finite-horizon contraction theorem by
iterated expectation. It would let Forge route an approximate relational goal
to exact transport duality rather than enumerate candidate algebraic
inequalities. The first implementation gate should be the cost certificate
and finite composition theorem, before any claim about mixing times.
